% !TEX root = ../print.tex
% Draft \S8 --- The admissible cone: the Gaussian ray and the Cin rays
%
% Source: draft/spatial-hemigroup-scale-space.md, \S8 (lines 479-520).
%
% CHANGED (lem:cin-rays): the draft states the expansion at infinity as
% Cin(z) = log z + gamma_E + O(1/z), with a check asking for a table anchor for gamma_E.
% The node states instead that there IS a finite constant C with Cin(z) = log z + C + O(1/z),
% which is proved here in three lines from the convergence of the cosine integral, and names
% the value C = gamma_E in a following remark with its citation. Nothing in the chapter, and
% nothing later, uses the value: prop:choquet-cone needs only the two growth rates. This
% keeps a [T] node from resting on an unproved numerical identity.
%
% prop:cin-origin-singularity is the one [A] node here. The draft leaves the origin
% behaviour of the Cin kernels as prose; it is a real claim resting on two theorems of Sato,
% and it supports prop:jet-regularity(4), so it is stated as a cited interface rather than
% asserted in a remark. Its ledger entry carries an unresolved reading, recorded there.

\chapter{The admissible cone: the Gaussian ray and the $\Cin$ rays}
\label{ch:cone}

\ref{thm:main-characterization} attaches to every admissible family a pair $(a,k)$, and
\ref{lem:admissible-cone} says the pairs form a convex cone that the exponents inherit
linearly. This chapter gives the cone its coordinates. The half-line cone had one
distinguished ray, the drift, sitting as a degenerate boundary; the line's cone has the
Gaussian ray, a genuine smoothing member, and it is extreme. The other extreme rays are the
exponents of the unit-step profiles.

% draft: Lemma 8.1
% twin: hcs lem:dickman-superposition (Hemigroup.dickman_superposition) --- same statement
% with Ein in place of Cin and the one-sided profile in place of the folded one.
\begin{lemma}[the $\Cin$ rays and the superposition]\label{lem:cin-rays}
\uses{lem:profile-integrability}
\lean{SpatialLine.cin_ray, SpatialLine.cin_elementary, SpatialLine.cin_expansion_zero,
SpatialLine.cin_expansion_top, SpatialLine.cin_superposition,
SpatialLine.cin_superposition_exists}\leanok
\notes{dickman-cin-rays}
Write $\Cin(z) := \int_0^z (1 - \cos v)\,v^{-1}dv$ for $z \ge 0$, extended evenly to $\RR$.
\begin{enumerate}
\item \emph{(The rays.)} For $\tau > 0$ the profile $k = \mathbf 1_{(0,\tau)}$ is admissible,
with Gaussian coefficient $0$, and its exponent is $\omega \mapsto \Cin(\tau\omega)$. The
function $\Cin$ is even and nondecreasing on $[0,\infty)$, with $\Cin(z) \le z^2/4$ and
$\Cin(z) = \tfrac14 z^2 + O(z^4)$ as $z \to 0$; and there is a finite constant $C$ with
$\Cin(z) = \log z + C + O(1/z)$ as $z \to \infty$.
\item \emph{(The superposition.)} Every admissible $F$, with Gaussian coefficient $a$ and
profile $k$, satisfies
\begin{equation}\label{eq:cone-superposition}
  F(\omega) = a\,\omega^2 + \int_{(0,\infty)} \Cin(\tau\omega)\,\varpi(d\tau),
  \qquad \omega \in \RR,
\end{equation}
where $\varpi$ is any positive measure on $(0,\infty)$ with $\varpi((x,\infty)) = k(x)$ for
almost every $x > 0$; such a $\varpi$ exists.
\end{enumerate}
\statusT\quad\emph{(Elementary throughout. The value $C = \gamma_E$ is
\ref{rem:cin-constant}; the article uses only that $C$ is finite, since what
\ref{prop:choquet-cone} needs is the two growth rates $z^2$ and $\log z$ and nothing else.)}
\end{lemma}

% CHANGED (proof of record 2026-09-09): rewritten to the machine-checked route, and two
% \uses edges removed with it. The printed proof reached admissibility through
% prop:fourier-toolbox(3) and the existence of varpi through lem:selfdecomposable-exponents.
% Neither is needed for what this lemma asserts. Admissibility here is exactly the two
% integrability conditions on the pair (a,k), which for a bounded compactly supported k are
% immediate; the passage to LE_s is lem:profile-integrability's own consequence clause and no
% clause of THIS node asks for it. And the tail measure is built by a quantile transform that
% consumes only "k nonincreasing" and "k vanishing at infinity" -- the second of which is
% FORCED by the integrability at infinity, not supplied by chapter 7. So the node is
% independent of chapters 2(3) and 7, and prop:choquet-cone inherits the shorter route.
% CHANGED (2026-09-11, draft sync, R101): the last paragraph said the tail measure of a BOUNDED
% profile puts infinite mass near the origin. It is the other way about: varpi((0,eps]) =
% k(0+) - k(eps), finite when k is bounded, so only an UNBOUNDED profile's tail measure has
% infinite mass there. Now "unbounded", as the draft already read. (The atom the quantile
% transform leaves at the origin when k is bounded, paragraph (2), is a different object and is
% rightly "bounded".)
\begin{proof}
\leanok
(1) Evenness of $\Cin$ needs no stipulation: the integrand $v \mapsto (1 - \cos v)/v$ is odd,
so the integral from $0$ is even, and $\Cin(-z) = \Cin(z)$. Monotonicity on $[0,\infty)$ is the
positivity of that integrand there. From $1 - \cos v \le v^2/2$ we get
$\Cin(z) \le \int_0^z v/2\,dv = z^2/4$ for $z \ge 0$, and for $z < 0$ by evenness, so the bound
holds on $\RR$. For the expansion at the origin, $|\cos v - (1 - v^2/2)| \le \tfrac{5}{96}v^4$
for $|v| \le 1$ gives $|(1-\cos v)/v - v/2| \le \tfrac{5}{96}v^3$ on $(0,1]$, and integrating
over $[0,z]$,
\[
  \Bigl|\Cin(z) - \tfrac14 z^2\Bigr| \ \le\ \tfrac{5}{96}\,z^4,
  \qquad |z| \le 1,
\]
which is the asserted $\Cin(z) = \tfrac14z^2 + O(z^4)$. The constant $\tfrac5{96}$ is not sharp:
the series gives $\Cin(z) = \tfrac14z^2 - \tfrac1{96}z^4 + O(z^6)$, so $\tfrac1{96}$ is the best
constant in the display, and the elementary $|\cos v - (1 - v^2/2)| \le v^4/24$ integrates to it.
It is kept because it is the constant the machine-checked proof carries, which it takes from the
library's cosine bound, and because nothing reads the value: what the statement asserts is
$O(z^4)$. For the expansion at infinity, split at
$v = 1$:
\[
  \Cin(z) = \Cin(1) + \log z - \int_1^z \frac{\cos v}{v}\,dv , \qquad z \ge 1,
\]
using $(1-\cos v)/v = v^{-1} - v^{-1}\cos v$ on $[1,z]$; and integrating by parts against the
primitive $\sin v/v$,
$\int_1^z v^{-1}\cos v\,dv = \sin z/z - \sin 1 + \int_1^z v^{-2}\sin v\,dv$. The last integrand
is dominated by $v^{-2}$, so $\int_1^\infty v^{-2}\sin v\,dv$ converges absolutely and, with
\[
  C := \Cin(1) + \sin 1 - \int_1^\infty \frac{\sin v}{v^2}\,dv ,
\]
one gets $\Cin(z) - \log z - C = -\sin z/z + \int_z^\infty v^{-2}\sin v\,dv$, whence
$|\Cin(z) - \log z - C| \le 2/z$ for $z \ge 1$. That is the expansion, with an explicit
constant in the error.

For the ray, $k = \mathbf 1_{(0,\tau)}$ is nonnegative and nonincreasing on $(0,\infty)$ and
satisfies both integrability conditions of \ref{lem:profile-integrability} at once, being
bounded by $1$ and supported in $(0,\tau)$: $\int_0^1 x\,k(x)\,dx \le \tfrac12$ and
$\int_1^\infty k(x)\,x^{-1}dx \le \tau$. So $(0,k)$ is a pair of the form \ref{eq:sd-profile},
and its exponent is
$\int_0^\tau (1-\cos\omega x)\,x^{-1}dx = \Cin(\tau\omega)$ by the substitution $v = \omega x$,
valid for every real $\omega$ because both sides vanish at $\omega = 0$ and $\Cin$ is even.

(2) A positive measure $\varpi$ on $(0,\infty)$ with $\varpi((x,\infty)) = k(x)$ for almost
every $x > 0$ is the image of Lebesgue measure on $(0,\infty)$ under the generalised inverse
$y \mapsto \sup\{u > 0 : k(u) > y\}$, restricted to $(0,\infty)$ to discard the atom the
transform leaves at the origin when $k$ is bounded. Only two properties of $k$ are used: it is
nonincreasing, and it tends to $0$ at infinity --- and the second is forced rather than
assumed, since $k$ nonincreasing and bounded below by a positive constant would make
$\int_1^\infty k(x)\,x^{-1}dx$ diverge like the harmonic integral. What comes out is the
sandwich $k(x{+}) \le \varpi((x,\infty)) \le k(x)$, whose two sides agree off the countably
many discontinuities of $k$; that is why the tail identity is asserted almost everywhere and
not everywhere, and a right-continuous version of $k$ is neither available nor needed.

Given such a $\varpi$, \ref{eq:cone-superposition} is the layer-cake exchange
\[
  \int_{(0,\infty)} \Cin(\tau\omega)\,\varpi(d\tau)
  = \int_0^\infty \frac{1 - \cos\omega x}{x}\,\varpi((x,\infty))\,dx
  = \int_0^\infty (1 - \cos\omega x)\,k(x)\,\frac{dx}{x},
\]
in which $\Cin(\tau\omega)$ is read as $\int_0^\tau (1-\cos\omega x)x^{-1}dx$ and the exchange
is Tonelli, for which $\varpi$ is $\sigma$-finite: the tail measure of an unbounded profile puts
infinite mass near the origin, but its restriction to $[\varepsilon,\infty)$ is finite for every
$\varepsilon > 0$. Adding $a\,\omega^2$ gives
\ref{eq:cone-superposition}.
\end{proof}

\begin{remark}[the constant]\label{rem:cin-constant}
The constant of \ref{lem:cin-rays}(1) is Euler's constant: $\Cin(z) = \gamma_E + \log z -
\mathrm{Ci}(z)$ with $\mathrm{Ci}$ the cosine integral and $\mathrm{Ci}(z) = O(1/z)$
(\sscite{dlmf2026}, (6.2.13) for the identity, and (6.2.20) with the expansions (6.12.3) and
(6.12.4) for the rate). The identification is not used anywhere in this article, which is
why \ref{lem:cin-rays} asserts only that the constant is finite.
\end{remark}

% CHANGED (proof of record 2026-09-09): the first two paragraphs of the proof are rewritten to
% the machine-checked route, which is where the domain argument changed. The printed proof asked
% for a comparison of Cin(z) with z^2 near 0 and with log z at infinity "uniformly comparable
% across omega on compact subsets of R minus {0}"; the checked route needs no uniform two-sided
% comparison at all. It needs two LOWER bounds -- Cin(z) >= (19/96)z^2 on [0,1] and
% Cin(z) >= (1 + log z)/2 past a threshold, both one line from the expansions of lem:cin-rays(1)
% -- and, on the bounded stretch between them, monotonicity of Cin alone. The Tonelli identities
% that turn the domain condition into the two integrability conditions of
% lem:profile-integrability are also written out now, because they are what the formalisation
% proves and what the review's "the domain condition is exactly right" was asserting.
% The existence of varpi is credited to lem:cin-rays(2) rather than to
% lem:selfdecomposable-exponents, which is where it is actually proved.
%
% CHANGED (proof of record 2026-09-09, wave 3): the last two paragraphs are rewritten to the
% machine-checked route, and lem:selfdecomposable-exponents is removed from \uses, an edge wave
% 2's rewrite had already replaced in the text and left in the list.
% (a) INJECTIVITY. The printed proof closed "hence the tails of varpi at almost every point,
% hence varpi", eliding two steps that turned out to be the price of the clause. Both are now
% written: the passage from tails equal a.e. to tails equal everywhere (a squeeze between two
% tails at points of the conull good set, then continuity of a measure from below), and the
% passage from tails to the measure. The second is NOT the pi-system argument the natural first
% draft reaches for and the statement annotation anticipated: the Choquet measure of an unbounded
% profile puts infinite mass near the origin, so the line carries no exhausting sequence of
% finite-measure half-lines, and the exhaustion has to be made after restricting to
% (epsilon, infinity), where the tails are finite. The paragraph says so.
% (b) EXTREMALITY. The printed proof identified the extreme rays of the domain and transported
% them, which is one argument; the checked development is three, and the two positive clauses do
% not use the transport at all. They are read at the profile: the summands of the Gaussian
% exponent have null \Levy{} measure, and the summands of a Cin exponent are a.e. constant on
% (0,tau) because each is at once nonincreasing (being a profile) and nondecreasing (being a
% constant minus a nonincreasing function). Only the classification uses the coordinates, and
% what it needs there is not "a positive measure that is not a multiple of a Dirac splits" as a
% piece of measure theory but an infimum argument, which is now written out. The rays and their
% folded \Levy{} densities, the printed paragraph's last sentence, are kept verbatim.
% CHANGED (2026-09-11, draft sync, R101): "the Choquet measure of a bounded profile has infinite
% mass near the origin", in the injectivity paragraph and in (a) above, is corrected to
% "unbounded": varpi((0,eps]) = k(0+) - k(eps) is finite when k is bounded, so a bounded profile
% does admit an exhausting sequence of finite-measure half-lines, and it is the unbounded case
% that forces the detour through (eps,infinity). The draft already read "unbounded".
% draft: Proposition 8.2
% twin: hcs prop:extreme-rays --- verbatim port, with the Gaussian ray in the drift's place.
% CHANGED (2026-09-19, article mirror of module B, R174): the closing clause glosses "Gaussian
% jitter" (the register pass's item 3: the term is used in the statement and introduced nowhere in
% module B, chapter 7's prose being the line paper's). A gloss of three words; no change of content.
\begin{proposition}[Choquet structure of the admissible cone]\label{prop:choquet-cone}
\uses{lem:cin-rays,lem:admissible-cone,thm:main-characterization,prop:fourier-toolbox,lem:profile-integrability,lem:quadratic-growth}
\lean{SpatialLine.choquet_cone_forward, SpatialLine.choquet_cone_surjective,
SpatialLine.choquet_cone_domain, SpatialLine.choquet_cone_linear,
SpatialLine.choquet_cone_injective,
SpatialLine.choquet_cone_extreme_gaussian, SpatialLine.choquet_cone_extreme_cin,
SpatialLine.choquet_cone_extreme_only}\leanok
\notes{dickman-cin-rays}
The superposition map of \ref{lem:cin-rays},
\[
  (a,\varpi) \ \mapsto\ F, \qquad
  F(\omega) = a\,\omega^2 + \int_{(0,\infty)} \Cin(\tau\omega)\,\varpi(d\tau),
\]
is a linear bijection from
$\bigl\{(a,\varpi) : a \ge 0,\ \varpi \ \text{a positive Borel measure on } (0,\infty),\
\int [\tau^2 \wedge (1 + \log_+\tau)]\,\varpi(d\tau) < \infty\bigr\}$
onto the cone of \ref{lem:admissible-cone}, with inverse $\varpi = -dk$, $k(x) =
\varpi((x,\infty))$, and $a$ the Gaussian coefficient. Consequently the extreme rays of the
cone are the Gaussian ray $\omega \mapsto \omega^2$ together with the one-parameter family
$\Cin(\tau\,\cdot)$, $\tau > 0$, whose folded \Levy{} densities are
$x^{-1}\mathbf 1_{(0,\tau)}(x)$; and every admissible family is a unique superposition of
$\Cin$ rays and Gaussian jitter, the term $a\,\omega^2$.
\statusT\quad\emph{(\textbf{Formalised} (2026-09-09), in eight declarations. The map, its
surjectivity with the inverse named here, the domain condition and the linearity rest on Lean
core alone. The injectivity of the map and the three extremality clauses rest, besides, on
uniqueness of the \Levy{} pair, \ref{prop:fourier-toolbox}(3) --- the one cited interface this
proposition spends, and the same one \ref{thm:matern} spends; no axiom was added for this
chapter. The domain condition is verified rather than asserted: it is
exactly the two integrability conditions of \ref{lem:profile-integrability} on the tail
profile, in both directions. Two steps of the injectivity argument are measure theory this
development had not needed before and are stated for arbitrary measures: an equality of tails
at almost every positive point is an equality at every positive point, and two measures carried
by $(0,\infty)$ with equal tails are equal.)}
\end{proposition}

\begin{proof}
\leanok
An admissible $F$ has, by \ref{thm:main-characterization}, a pair $(a,k)$ of the form
\ref{eq:sd-profile}; a positive measure $\varpi$ on $(0,\infty)$ with
$k(x) = \varpi((x,\infty))$ almost everywhere exists by \ref{lem:cin-rays}(2), and
\ref{eq:cone-superposition} is then the displayed representation. Conversely a pair
$(a,\varpi)$ in the stated domain gives back an admissible $F$: its profile is the tail
$k(x) = \varpi((x,\infty))$, nonnegative and nonincreasing, and the two integrability
conditions of \ref{lem:profile-integrability} are the domain condition read through Tonelli,
\[
  \int_0^1 x\,\varpi((x,\infty))\,dx = \int \tfrac12(\tau \wedge 1)^2\,\varpi(d\tau),
  \qquad
  \int_1^\infty \varpi((x,\infty))\,\frac{dx}{x} = \int \log_+\tau\,\varpi(d\tau),
\]
two more instances of the layer cake of \ref{lem:cin-rays}(2), with the weights $x$ cut off
above $1$ and $x^{-1}$ cut off below $1$. Since
$\tfrac12(\tau\wedge1)^2 + \log_+\tau \le \tau^2 \wedge (1 + \log_+\tau)
 \le 2\bigl[\tfrac12(\tau\wedge1)^2 + \log_+\tau\bigr]$ for $\tau > 0$, the domain condition and
the two integrability conditions are the same condition. So the map is onto the cone, and the
domain is what the statement says.

That the superposition integral is finite at one $\omega \ne 0$ if and only if at all, and if
and only if the domain condition holds, follows. Sufficiency is the paragraph just given
together with \ref{lem:quadratic-growth}, which makes an admissible exponent finite at every
frequency. For necessity it is enough to dominate $\tau^2 \wedge (1 + \log_+\tau)$ by a
multiple of $\Cin(\tau|\omega|)$ pointwise on $(0,\infty)$, and three regimes do that, using
only \emph{lower} bounds on $\Cin$. Below $\tau|\omega| = 1$ the expansion at the origin gives
$\Cin(z) \ge \tfrac{19}{96}z^2$, which dominates $\tau^2$ up to the constant
$|\omega|^{-2}$. Past a threshold the expansion at infinity gives
$\Cin(z) \ge \tfrac12(1 + \log z)$, which dominates $1 + \log_+\tau$ up to a constant once
$\tau$ is large enough that $\log(\tau|\omega|)$ is at least half of $\log\tau$. On the bounded
stretch between the two, monotonicity alone suffices:
$\Cin(\tau|\omega|) \ge \Cin(1) \ge \tfrac{19}{96}$ there, while $1 + \log_+\tau$ is bounded.
The map is linear in $(a,\varpi)$ by construction, both clauses being additivity of the
integral in the measure and positive homogeneity. It is injective because the pair of $F$ is
unique (\ref{prop:fourier-toolbox}(3)): the pair determines $a$ and the \Levy{} measure
$k(x)x^{-1}dx$, hence $k$ up to a null set, hence the tails $\varpi((x,\infty))$ at almost
every $x > 0$. Two steps carry that to $\varpi$ itself. First, the exceptional set being null,
every interval $(x, x + \varepsilon)$ contains a point $y$ at which the two tails agree, and
$\varpi_1((x + \varepsilon,\infty)) \le \varpi_1((y,\infty)) = \varpi_2((y,\infty)) \le
\varpi_2((x,\infty))$; letting $\varepsilon$ decrease and using continuity of a measure along
the increasing sets $(x + \varepsilon, \infty)$ gives equality of the tails at \emph{every}
$x > 0$. Second, the tails at positive points are finite, by the domain condition, so for each
$\varepsilon > 0$ the restrictions of $\varpi_1$ and $\varpi_2$ to $(\varepsilon,\infty)$ are
finite measures agreeing on every half-line $(-\infty,t]$, since $\varpi((\varepsilon,t])$ is a
difference of two finite tails; two finite measures on $\RR$ agreeing on all such half-lines are
equal. Letting $\varepsilon$ decrease again, and using that $\varpi_1$ and $\varpi_2$ give no
mass to $(-\infty,0]$, gives $\varpi_1 = \varpi_2$. The detour through a finite restriction is
not avoidable by a direct $\pi$-system argument on the line: the Choquet measure of an unbounded
profile has infinite mass near the origin, so no exhausting sequence of finite-measure
half-lines exists before the restriction is made.

Extremality of the two named families is read at the profile, not through the coordinates. If
$\omega^2 = F_1 + F_2$ with both summands admissible, then $F_1 + F_2$ has pair $(1,0)$ by
uniqueness, so $k_1 + k_2$ vanishes almost everywhere on $(0,\infty)$, so each $k_i$ does, and
$F_i(\omega) = a_i\,\omega^2$ with $a_i \ge 0$: the Gaussian ray is extreme. If
$\Cin(\tau\,\cdot) = F_1 + F_2$, then uniqueness gives $a_1 = a_2 = 0$ and
$k_1 + k_2 = \mathbf 1_{(0,\tau)}$ almost everywhere; on the conull set inside $(0,\tau)$ where
that holds, each $k_i$ is nonincreasing, being a profile, and nondecreasing, being what is left
of a constant after subtracting a nonincreasing function, hence constant there, while above
$\tau$ both vanish by nonnegativity. So $k_i = c_i\,\mathbf 1_{(0,\tau)}$ almost everywhere and
$F_i = c_i\,\Cin(\tau\,\cdot)$ with $c_i \ge 0$: each $\Cin$ ray is extreme.

That there are no other extreme rays is where the coordinates are used, and where the
injectivity just proved is spent. Let $F$ be extreme with pair $(a,\varpi)$. If $\varpi = 0$
then $k$ vanishes almost everywhere and $F(\omega) = a\,\omega^2$. If $\varpi \ne 0$ then
$a = 0$: otherwise splitting $F$ as $a\,\omega^2$ plus the pure-jump part makes $a\,\omega^2$
equal to $c_1F$ for some $c_1$, and either $c_1 = 0$, forcing $a = 0$, or $F$ is a multiple of
$\omega^2$, which by uniqueness forces $\varpi = 0$. Now suppose some $s > 0$ has
$\varpi((0,s]) > 0$ and $\varpi((s,\infty)) > 0$. The restrictions of $\varpi$ to $(0,s]$ and to
$(s,\infty)$ split $F$ into two admissible exponents; extremality makes the first of them
$c_1F$; injectivity turns that into an equality of measures, the restriction of $\varpi$ to
$(0,s]$ against $c_1\varpi$; evaluating both sides on $(s,\infty)$ forces $c_1 = 0$, and then
$\varpi((0,s]) = 0$, a contradiction. So every $s$ has mass on at most one side. Put
$\tau := \inf\{s > 0 : \varpi((0,s]) > 0\}$, a set that is nonempty because $\varpi \ne 0$.
Then $\varpi((x,\infty)) = 0$ for $x > \tau$, by the defining property of the infimum together
with the hypothesis just established, hence also at $x = \tau$; and $\varpi((0,x]) = 0$ for
$0 < x < \tau$, by minimality. So $\varpi$ is carried by the single point $\tau$, which is
positive since otherwise $\varpi$ would be $0$, and $F = c\,\Cin(\tau\,\cdot)$ with
$c = \varpi(\{\tau\})$, finite because $\varpi((\tau/2,\infty))$ is. This is the map's own
picture: $(1,0) \mapsto \omega^2$, the Gaussian ray, and
$(0,\delta_\tau) \mapsto \Cin(\tau\,\cdot)$, whose profile is
$k(x) = \delta_\tau((x,\infty)) = \mathbf 1_{(0,\tau)}(x)$ and whose folded \Levy{} density is
therefore $x^{-1}\mathbf 1_{(0,\tau)}(x)$.
\end{proof}

% CHANGED (2026-09-19, article mirror of module B, R174): three register repairs inside the
% remark, from the module B register pass's item 2 (notes/reviews/register-b/README.md).
% (a) "none of them extreme" is FALSE at one point of one family and now carries the exception:
% the stable family ends on the Gaussian ray at alpha = 2, which prop:choquet-cone makes extreme.
% Safe direction: the remark claims less. (b) The chain ID_s > SD_s > GGC_s was printed with no
% gloss -- the three symbols occur nowhere else in the blueprint and are defined in macros.tex
% only -- and the three classes are now named where they are used. (c) "which rests on a cited
% interface and is not yet machine-checked" is stale since R160: prop:student-t(3) is
% machine-checked under its hypothesis, and what is cited is that hypothesis (rem:student-ggc,
% ledger A18, unadmitted). The sentence now says which.
% CHANGED (2026-09-11, draft sync, R107): "all three lie in the Thorin subcone
% (prop:thorin-subclass)" cited the memberships to a node that, since R79, states none of them:
% its declarations prove the representation, the bridge identity and the strict inclusion. Each
% membership is now cited to the node that states it, as R79's paragraph after
% prop:thorin-subclass's proof does, the Student-t one being prop:student-t(3), cited and not
% machine-checked. (The ledger entry is named in prose, not by macro, so that this remark does
% not project it as a source.) Mirrored at draft Remark 8.3.
\begin{remark}[what the cone looks like, and what changed]\label{rem:cone-shape}
In the log-displacement coordinate of \ref{rem:displacement-dictionary} the proposition reads:
every nonincreasing profile is a unique positive stack of unit steps, and every admissible
family is that stack plus one number $a$, a point \emph{off} the curve. The three distinguished
families of this article are slices cut by different principles: the
symmetric stable family by dilation symmetry (\ref{cor:semigroup-case}), the \Matern{} family as
the single-atom member of the Thorin subclass (\ref{thm:matern}), and the Student-t family as
the image of the causal Bessel corner under the bridge (\ref{prop:student-t}). No member of the
three is an extreme ray of the cone, with one exception at an endpoint: the stable family closes
at $\alpha = 2$ on the Gaussian ray (\ref{prop:stable-family}), which is extreme. All three
lie in the Thorin subcone, each by the node that states it: the \Matern{} family as the one-atom
case of \ref{prop:thorin-subclass}(2), the stable family by its Thorin density
(\ref{prop:stable-family}), and the Student-t family by \ref{prop:student-t}(3), whose Thorin
half is proved under the hypothesis that the causal delay law is a generalized gamma
convolution, a cited fact (\ref{rem:student-ggc}). The $\Cin$ rays are the paradigm of
what lies outside it: profiles with a hard cutoff, whose kernels carry derivative singularities
at the displacements $\pm\tau, \pm2\tau, \dots$, echoes of the step
(\ref{lem:cin-delay-equation}). Two things are new against the causal cone. The Gaussian ray is
an extreme ray of every subcone in the chain
$\IDs \supset \SDs \supset \GGCs$ --- the symmetric infinitely divisible laws, the symmetric
self-decomposable ones, and the symmetric Thorin subclass of \ref{prop:thorin-subclass} --- being present
in each with $k \equiv 0$, whereas the corresponding causal ray, the drift, was the degenerate
member. And the $\Cin$ rays, unlike their causal counterparts, are not the image of anything
under the bridge of Chapter~\ref{ch:bridge} (\ref{prop:bridge-strictness}). The tail statement
of the causal remark holds here in a different shape but with the same conclusion: by
\ref{prop:moments-tails}(2) a $\Cin$ kernel has tails of order $e^{-|x|\log|x|/\tau}$, lighter
than any exponential and heavier than any Gaussian; any admissible kernel with a jump catalogue
has tails at least that heavy, whatever its Gaussian coefficient; and the only admissible
kernels with Gaussian tails are the Gaussians themselves. No truncation of the Gaussian is
admissible.
\end{remark}

% draft: Lemma 8.4
% NEW: the exact twin of the delay equation of the Dickman function, with the average of the
% two neighbours in place of the one delayed value, because displacements come in both signs.
% twin: hcs has no node for the Dickman delay equation; the causal article states it in prose
% in its cone remark.
% CHANGED (proof 2026-09-10): the two Lean declarations named below gained the hypothesis that
% the density is measurable (AEMeasurable p volume). This is a REPAIR, not a convenience: both
% declarations were FALSE without it. The hypothesis mu = volume.withDensity (ofReal p) sees a
% non-measurable p only through the measurable simple functions below it, so raising p on a
% Bernstein set -- inner measure zero, full outer measure -- leaves the hypothesis intact while
% the conclusion fails on a set no null set contains; the derivative form fails likewise, at a
% test function supported where p was altered but p(.+-tau) was not, one side being junk-zero
% and the other not. The prose statement of record is unaffected: "the density of the symmetric
% law" is a density, and a density is measurable. Every call site supplies the hypothesis, the
% density being produced by prop:kernel-regularity. Wave 5 recorded the missing hypothesis as a
% review question; wave 6 found it to be a defect and repaired it.
\begin{lemma}[the $\Cin$ kernels satisfy a delay equation]\label{lem:cin-delay-equation}
\uses{lem:cin-rays,prop:kernel-regularity,prop:fourier-uniqueness,thm:increments-levy}
\lean{SpatialLine.cin_delay_equation, SpatialLine.cin_delay_equation_deriv}\leanok
\notes{dickman-cin-rays}
Let $p_\tau$ be the density of the symmetric law with exponent $\Cin(\tau\,\cdot)$ --- folded
profile $\mathbf 1_{(0,\tau)}$, two-sided \Levy{} density
$\tfrac12|x|^{-1}\mathbf 1_{(0,\tau)}(|x|)$, no Gaussian part --- and let $P_\tau$ be its
distribution function. Then, in the sense of distributions,
\[
  x\,p_\tau(x) = P_\tau(x) - \tfrac12\bigl[P_\tau(x-\tau) + P_\tau(x+\tau)\bigr],
  \qquad\text{equivalently}\qquad
  x\,p_\tau'(x) = -\tfrac12\bigl[p_\tau(x-\tau) + p_\tau(x+\tau)\bigr].
\]
\statusT\quad\emph{(The identity propagates whatever singularity $p_\tau$ has at the origin to
$\pm\tau$, thence to $\pm2\tau$, and so on, with one more order of smoothness at each step,
exactly as for the Dickman kernel; \ref{prop:cin-origin-singularity} says what the singularity
at the origin is. The equality is asserted in the sense of distributions, which is what the
transform argument delivers; both sides are locally integrable, so it is also an
almost-everywhere identity. \textbf{Machine-checked} (2026-09-10) in both forms, on Lean core.
Read at the law rather than at the density the two displays are one statement, an identity
between two finite signed measures, and that is how the proof below is organised; the passage
back to the density is a change of variables for the first display and one integration by parts
for the second. The two Lean declarations carry one hypothesis the display above leaves
implicit --- that the density is measurable --- because the Lean statements quantify over an
arbitrary $p$ with $\mu = \int p\,dx$, and such a hypothesis constrains $p$ only through the
measurable functions below it; the prose statement, which speaks of \emph{the} density, is
unaffected.)}
\end{lemma}

% CHANGED (proof of record 2026-09-10): rewritten to the machine-checked route. The printed
% proof convolved p with the SIGNED measure y nu_2(dy) = (1/2) sgn(y) 1_{|y|<tau} dy and then
% differentiated the resulting a.e. identity. The checked route keeps the two halves of that
% weight positive until the last step -- the first is Lebesgue measure on [0,tau), the second
% its reflection -- so what is convolved is a positive measure, its transform is the transform
% of a convolution, and no signed measure is convolved anywhere. Only at the end are two
% positive finite measures subtracted, and uniqueness is applied to the Jordan recombination.
% The differentiation of the a.e. identity, which the printed proof performs formally, is
% replaced by the pairing against a test function, where it is an integration by parts.
% CHANGED (2026-09-11, draft sync, R102): eq:cin-delay-measures defined nu_tau := nu(. - tau),
% which read as nu_tau(A) = nu(A - tau) is the image of nu under x -> x + tau: density
% mu((x-2tau, x-tau]). The next sentence gives nu_tau the density mu((x, x+tau]) = (density of
% nu)(x + tau), and the transform display gives it e^{-i omega tau} nu-hat under the kernel
% e^{i omega x} the proof computes with; both are the image under x -> x - tau, and so is the
% Lean (CinDelayForm.lean reads nu at x + tau). The density fixes the direction whatever the sign
% convention. Now "the image of nu under x -> x - tau", as the draft already read.
% CHANGED (2026-09-11, post-sync follow-up, R110): the transforms in this proof were computed
% with the kernel e^{+i omega x}, against def (2.0)'s e^{-i omega x} of prop:fourier-uniqueness,
% which the proof cites. Now computed with e^{-i omega x}: the moment transform is
% -i tau Cin'(tau omega) c(omega), the window transform (1 - e^{-i omega tau})/(i omega), the shift
% factor e^{+i omega tau}, and the difference -2i (1 - cos omega tau) c(omega)/omega. Every
% displayed quantity is the complex conjugate of the old one, c being real, so the identity
% eq:cin-delay-measures and the route are unchanged. The Lean (CinDelayForm.lean) computes with
% Mathlib's charFun, whose kernel is e^{+i omega x}; the two are conjugate and uniqueness holds
% for either, so this is a convention alignment, not a change of route.
\begin{proof}
\leanok
Write $\mu$ for the law, $c(\omega) = \int\cos(\omega x)\,\mu(dx) = e^{-\Cin(\tau\omega)}$ for
its transform, and $\kappa_\tau$ for Lebesgue measure on $[0,\tau)$. Everything follows from one
identity between finite signed measures on $\RR$:
\begin{equation}\label{eq:cin-delay-measures}
  2\,x\,\mu(dx) \;=\; \nu - \nu_\tau ,
  \qquad \nu := \mu * \kappa_\tau, \quad
  \nu_\tau := \text{the image of $\nu$ under } x \mapsto x - \tau .
\end{equation}
By Tonelli, $\nu$ has Lebesgue density $x \mapsto \mu((x-\tau,x])$ and $\nu_\tau$ has density
$x \mapsto \mu((x,x+\tau])$; both are finite measures of total mass $\tau$. The left-hand side
is finite because $\mu$ has a finite first absolute moment: $\Cin(z) \le z^2/4$, so the exponent
of $\mu$ is quadratically flat, the truncation inequality of \ref{thm:increments-levy} gives
$\mu\{|x| > R\} = O(R^{-2})$, and the layer cake integrates that tail.

\emph{The transforms.} Since $\mu$ is symmetric, $\int x\,e^{-i\omega x}\mu(dx)$ has no real
part, the integrand's cosine half being odd, so it equals $-i\int x\sin(\omega x)\,\mu(dx)$. The
first absolute moment lets $c$ be differentiated under the integral sign, giving
$c'(\omega) = -\int x\sin(\omega x)\,\mu(dx)$; and differentiating $c = e^{-\Cin(\tau\cdot)}$ by
the chain rule gives $c'(\omega) = -\tau\,\Cin'(\tau\omega)\,c(\omega)$ with
$\Cin'(v) = (1-\cos v)/v$. The two readings of $c'$ are the elementary delay ODE
\begin{equation}\label{eq:cin-delay-ode}
  \omega\,c'(\omega) = -\,c(\omega)\,\bigl(1 - \cos\tau\omega\bigr),
\end{equation}
and together they give
$\int x\,e^{-i\omega x}\mu(dx) = -i\,\tau\,\Cin'(\tau\omega)\,c(\omega)$, which for
$\omega \ne 0$ is $-i\,(1-\cos\tau\omega)\,c(\omega)/\omega$. On the other side, the transform of
a convolution is the product, so
\[
  \hat\nu(\omega) = c(\omega)\int_0^\tau e^{-i\omega u}\,du
    = c(\omega)\,\frac{1 - e^{-i\omega\tau}}{i\omega},
  \qquad \hat\nu_\tau(\omega) = e^{i\omega\tau}\,\hat\nu(\omega),
\]
whence, for $\omega \ne 0$,
\[
  \hat\nu(\omega) - \hat\nu_\tau(\omega)
    = c(\omega)\,\frac{(1 - e^{-i\omega\tau})(1 - e^{i\omega\tau})}{i\omega}
    = c(\omega)\,\frac{2 - 2\cos\omega\tau}{i\omega}
    = -2i\,\frac{1 - \cos\omega\tau}{\omega}\,c(\omega),
\]
which is twice the transform of $x\,\mu(dx)$, that is, the transform of the left-hand side of \ref{eq:cin-delay-measures}; at $\omega = 0$ both
sides vanish, the left by symmetry and the right because $\nu$ and $\nu_\tau$ have equal mass.

\emph{Uniqueness.} Both sides of \ref{eq:cin-delay-measures} are signed, so
\ref{prop:fourier-uniqueness} is applied after recombination: writing $x^{\pm}$ for the positive
and negative parts of $x$, the two \emph{positive} finite measures
$2\,x^{+}\mu(dx) + \nu_\tau$ and $2\,x^{-}\mu(dx) + \nu$ have the same transform
by the display above, hence are equal, which is \ref{eq:cin-delay-measures}.

\emph{The two displays.} The first is \ref{eq:cin-delay-measures} read at Lebesgue densities.
With $\mu(dx) = p_\tau(x)\,dx$ the left-hand side has density $2x\,p_\tau(x)$ and the right-hand
side has density
\[
  \mu((x-\tau,x]) - \mu((x,x+\tau])
  = 2\Bigl(P_\tau(x) - \tfrac12\bigl[P_\tau(x-\tau) + P_\tau(x+\tau)\bigr]\Bigr),
\]
and two integrable functions with the same integral against every bounded measurable test
function agree almost everywhere --- it is enough to test against the sign of their difference.
The second display is \ref{eq:cin-delay-measures} paired with $\varphi'$ for a test function
$\varphi \in C_c^\infty$. Since $\int_0^\tau \varphi'(x+u)\,du = \varphi(x+\tau) - \varphi(x)$,
and likewise for $\nu_\tau$ with $x$ shifted, pairing gives
\[
  2\int x\,\varphi'(x)\,\mu(dx)
   = \int\bigl[\varphi(x+\tau) + \varphi(x-\tau)\bigr]\mu(dx) - 2\int \varphi\,d\mu ,
\]
that is $\int (x\varphi)'\,d\mu = \tfrac12\int[\varphi(\cdot+\tau) + \varphi(\cdot-\tau)]\,d\mu$;
and one integration by parts against the density turns that into the second display. One route
to the same identity convolves $\mu$ with the \emph{signed} weight
$y\,\nu_2(dy) = \tfrac12\sgn(y)\mathbf 1_{|y|<\tau}\,dy$, which is $\tfrac12$ of $\kappa_\tau$
minus its reflection. The computation above is that one with the two halves kept positive until
the subtraction at the end, so that no signed measure is convolved anywhere.
\end{proof}

% NEW NODE (2026-09-15, module B step 2, R159; the postdoc's P-0008): the folding translation,
% made a node of its own and machine-checked. Until now it was asserted in prose in three places
% -- chapter 2's folding-convention paragraph, ledger A10's "what the entry does NOT carry"
% bullet, and the annotation of prop:cin-origin-singularity below -- and checked nowhere. It is
% [T], it is elementary, and it is the second of the two obstacles the wave-6 note records
% against admitting A10: admitting the three regimes in their folded form would put this
% translation behind an admitted name. With the lemma proved, an admission would carry the
% source's own statement and nothing more. Nothing in the article claimed less before: the
% translation was always asserted, and it is now checked.
% CHANGED 2026-09-19 (R180, second external review round, referee A findings 4 and 5): THE
% ANNOTATION ONLY -- no clause, no tag, no status. The referee found two hypotheses of the cited
% theorems that this lemma verifies nowhere: Sato's (53.14) driftless finite-variation form with
% its (53.15), and the left-continuity normalisation of his k-function (p. 403, repeated in the
% hypothesis of Thm. 28.4 p. 191). Both hold here -- the first from k <= c under the hypothesis
% c < infinity, the drift from the cosine-transform specification of a symmetric law and the
% Gaussian part from a = 0; the second by a modification on a countable set, under which the
% \Levy{} measure, the exponent, the two one-sided limits and both comparison functions are
% unchanged -- and the whole point of this lemma is that such matchings are made on the page, so
% they are written into the annotation rather than left to a reader. NOT into the statement: the
% lemma is \leanok and folding_translation proves the four conjuncts it has; an added conclusion
% would be an unproved clause under a \lean tag. The left-continuity step is a step the three A10
% names carry beyond their pages and is added to the per-name list at the head of
% trust-boundary.txt; the finite-variation one is not, being implied by the hypotheses the
% specifications already carry.
\begin{lemma}[the folding translation]\label{lem:folding-translation}
\uses{lem:profile-integrability}
\lean{SpatialLine.folding_translation}\leanok
Let $(a,k)$ be a pair of the form \ref{eq:sd-profile} with $a = 0$, and suppose the folded
profile has a finite right limit $c := k(0{+})$ at the origin. Put $k_2(x) := \tfrac12 k(|x|)$
for $x \in \RR$. Then $k_2$ is a two-sided profile in the sense of the probabilistic sources:
it is nonnegative, nonincreasing on $(0,\infty)$, nondecreasing on $(-\infty,0)$, and
$\int_\RR (1 \wedge x^2)\,k_2(x)\,|x|^{-1}dx < \infty$. Its \Levy{} measure
$k_2(x)|x|^{-1}dx$ on $\RR \setminus \{0\}$ is the unfolding of $k(x)x^{-1}dx$:
\[
  \int_\RR \bigl(1 - \cos\omega x\bigr)\,\frac{k_2(x)}{|x|}\,dx
   = \int_0^\infty \bigl(1 - \cos\omega x\bigr)\,\frac{k(x)}{x}\,dx
   \qquad (\omega \in \RR),
\]
so the two forms specify the same symmetric law. Moreover $k_2(x) + k_2(-x) = k(x)$ for every
$x > 0$, and both one-sided limits at the origin exist, with
\[
  k_2(0{+}) = k_2(0{-}) = \tfrac12 c,
  \qquad\text{so that}\qquad
  k_2(0{+}) + k_2(0{-}) = c \quad\text{and}\quad k_2(0{+}) - k_2(0{-}) = 0 .
\]
\statusT\quad\emph{(Machine-checked in full and on Lean core. The lemma is the dictionary
between this article's folded convention and the two-sided convention of the probabilistic
sources, and it is stated because the sources state their results about a self-decomposable
density at the origin in terms of the two constants $k_2(0{+}) + k_2(0{-})$ and
$k_2(0{+}) - k_2(0{-})$: the first is the folded $k(0{+})$ and the second vanishes. The
hypothesis that $c$ is finite is used only for the two limits; the structural conditions and
the exponent identity hold with no hypothesis at the origin. The two-sided profile is a
vocabulary for citations and not a class this development quantifies over: it is defined in
\texttt{SpatialLine/TwoSidedDefs.lean} and appears nowhere else in the library.
\textbf{Two further hypotheses of the sources, and why $k_2$ meets them} (2026-09-19, R180; the
second external review found both unmatched). The singular and threshold regimes are stated for a
law of the source's driftless finite-variation form --- Sato's (53.14), which asks
$\int_{|x|<1}k_2(x)\,dx < \infty$ as its (53.15) and has no drift term and no Gaussian part. Under
the hypothesis $c < \infty$ carried above, $k_2 \le \tfrac12c$ on each side of the origin by
monotonicity, so $\int_{|x|<1}k_2(x)\,dx \le c < \infty$; the drift vanishes because the law is
specified through its cosine transform, which fixes the symmetrisation and leaves no first-order
term; and there is no Gaussian part because $a = 0$. That is why the finiteness of $c$ is the
hypothesis those two regimes carry and the smooth one does not need in the same way.
\emph{Sato's $k$-function is moreover the version left-continuous on $(0,\infty)$ and
right-continuous on $(-\infty,0)$}, and $k_2$ as defined here need not be either. Nothing is lost,
and three invariances say so: passing to $x \mapsto k_2(x{-})$ on $(0,\infty)$ and $k_2(x{+})$ on
$(-\infty,0)$ changes $k_2$ on a countable set, hence changes neither the \Levy{} measure
$k_2(x)|x|^{-1}dx$ nor the exponent nor the law; it changes neither one-sided limit at the origin,
so $c$ and the antisymmetric constant are the same; and it changes neither $K$ nor $L$, both being
integrals of $k_2$. So the source's theorems, applied to the modification, conclude about the same
law and with the same comparison functions. The Lean specifications of
\texttt{SpatialLine/TwoSidedDefs.lean} quantify over a monotone $k$ without the continuity
normalisation, and that is the step: it is recorded at the three names in
\texttt{blueprint/trust-boundary.txt}.)}
\end{lemma}

\begin{proof}
\leanok
The three structural conditions are the corresponding conditions on $k$ read through
$x \mapsto |x|$: nonnegativity is immediate, and for $x \le y < 0$ one has $|x| \ge |y| > 0$, so
$k_2(x) \le k_2(y)$ because $k$ is nonincreasing on $(0,\infty)$. For the two integrals, note
first that the map $x \mapsto -x$ preserves Lebesgue measure and that both integrands below are
even, $\cos$ being even and $k_2$ being a function of $|x|$; so each integral over $\RR$ is
twice the corresponding integral over $(0,\infty)$, where $k_2 = \tfrac12 k$. The two factors
$2$ cancel, which gives the exponent identity, and the same computation turns
$\int_\RR (1 \wedge x^2)k_2(x)|x|^{-1}dx$ into $\int_0^\infty (1 \wedge x^2)k(x)x^{-1}dx$, which
is finite: on $(0,1)$ it is $\int_0^1 x\,k(x)\,dx$ and on $(1,\infty)$ it is
$\int_1^\infty k(x)x^{-1}dx$, the two conditions of \ref{lem:profile-integrability}.

The pointwise identity is $k_2(x) + k_2(-x) = \tfrac12 k(x) + \tfrac12 k(x)$ for $x > 0$. For the
limits, $k_2$ agrees with $\tfrac12 k$ on $(0,\infty)$, so its right limit is $\tfrac12 c$; and
$k_2(x) = \tfrac12 k(-x)$ on $(-\infty,0)$, where $x \mapsto -x$ carries the left neighbourhood
filter at the origin to the right one, so its left limit is $\tfrac12 c$ as well. The sum and
the difference follow.
\end{proof}

% additive as a NODE, prose in the draft (\S8, the paragraph after Lemma 8.4). It is a real
% claim, it rests on two theorems of Sato, and prop:jet-regularity(4) depends on it, so it is
% stated as a cited interface rather than asserted in a remark.
% CHECK (draft, checks for \S8, item 5): the constant in Sato's (53.30) and the reading of
% its cosine factor are not yet page-verified; the ledger entry says so, and the node claims
% only the order of the singularity, not the constant.
% CHANGED (proof 2026-09-09, sweep): the Lean declarations named below gained the hypothesis
% that the kernels are SYMMETRIC. This is the review's R7 defect, swept for across chapters
% 5-13 by wave 2's merge: a declaration whose hypotheses constrain a measure only through its
% COSINE transform -- which fixes the symmetrisation and nothing more -- may not conclude
% anything about the measure itself. Symmetry is what (A3) gives every family of this article,
% so the hypothesis narrows nothing; the statement of record below is unaffected, the gap being
% in the LEAN quantification and not in the prose.
% Affected: cin_origin_singularity_unbounded, _smooth and _threshold, each of which asserts
% that the kernel HAS A DENSITY with a prescribed two-sided bound or smoothness. Counterexample
% without it: replace the symmetric law of density q by the probability measure of density
% 2q(x)1_{x>0}, which has the same cosine transform and density 0 on the negative half-line.
% The fourth declaration, SpatialLine.cin_origin_singularity_cin_ray, mentions no measure and
% is unaffected.
% CHANGED (author's decision 2026-09-10): the \uses edge to lem:cin-delay-equation is REMOVED.
% It was vestigial: it recorded the propagation clause, which review Q8 moved out of this
% statement into rem:cin-propagation on 2026-09-09, and wave 6 established that what remains --
% the three regimes and the Cin-ray sentence -- never consumed the delay equation, the regimes
% quantifying over an arbitrary admissible F with Gaussian coefficient 0 and speaking only of
% its kernel near the origin. The dependency graph is a statement of record, so an edge that is
% merely harmless is still false. The annotation's \ref to lem:cin-delay-equation stays: it says
% what the node does NOT depend on, and rem:cin-propagation, which does depend on it, carries
% the reference that matters. A \ref is not a \uses.
% CHANGED (2026-09-11, draft sync, R100): the threshold regime read "phi_1 is continuous at no
% better than logarithmic order, phi_1(x) ~ L(x)", which contradicts its own comparison: k is
% nonincreasing with k(0+) = c = 1, so k <= 1, the integrand (1 - k(u))/u is nonnegative, K >= 1
% on (0,1), and L(x) >= log(1/|x|) -> infinity. A function comparable to L is unbounded near the
% origin and does not extend continuously there. Now "is unbounded, phi_1(x) ~ L(x)". The word
% adds no claim: unboundedness is a consequence of the two-sided comparison, which is all that
% Skeleton.cin_origin_singularity_threshold states (its lower bound c_1 * L with c_1 > 0), exactly
% as the Lean docstring at the c < 1 regime says of "unbounded" there. Mirrored at draft
% Proposition 8.5, which had dropped the phrase.
% CHANGED (2026-09-15, module B step 2, R159): the ANNOTATION only; the statement, the tag and
% the status are untouched and the node stays \notready. The closing sentences said that two
% things stood in the way of admitting ledger A10 -- that nothing spends the three regimes, and
% that admitting them in their folded form would put the folding translation behind an admitted
% name -- and both are now answered: lem:folding-translation above is the translation, checked
% on Lean core; SpatialLine.cin_origin_singularity_unbounded_of_sato, _smooth_of_sato and
% _threshold_of_sato derive these three statements from the source's own sentences taken as
% hypotheses, also on Lean core; and cor:origin-boundedness below reads all three regimes. The
% node remains \notready because the interface is still unadmitted, which is what \notready
% means here.
% CHANGED (2026-09-15, module B step 2, R162): the node is \leanok, THE AUTHOR HAVING ADMITTED
% ledger A10. The STATEMENT IS UNCHANGED, word for word; what changed is the tag and the status.
% The three regimes are now SpatialLine.cin_origin_singularity_unbounded, _smooth and _threshold,
% proved in SpatialLine/OriginSingularity.lean -- each one application of the corresponding
% _of_sato theorem to the corresponding axiom -- where they were Skeleton declarations carrying
% sorry, and those are deleted. The three axioms are SpatialLine.sato_origin_singular,
% sato_origin_smooth and sato_origin_threshold, declared in SpatialLine/Interfaces.lean and
% stated AT SATO'S LETTER, over the source's own two-sided profile and its two one-sided limits;
% the translation into this article's folded convention is lem:folding-translation above and is
% NOT carried by them. A proof environment is added below for the reason R157 added one at
% prop:kernel-regularity: a \leanok [A] node that prints nothing between its citation and its
% conclusion leaves this article's own step -- here the passage between the two conventions -- in
% the annotation alone. The \uses edge to lem:folding-translation is added, that lemma being what
% the three proofs now spend. Neither headline theorem is affected: neither reaches this chapter,
% and both #print axioms blocks are byte-identical across the admission.
% CHANGED (2026-09-19, external review of module B, R171): THE SINGULAR REGIME WAS FALSE, AND SO
% WAS THE ADMITTED AXIOM. The regime read "phi_1(x) asymp |x|^{c-1}" for c < 1, dropping the
% slowly varying factor K of Sato's (53.25) -- the same K the threshold clause already carried
% inside its L. The referee's counterexample: the folded profile k(x) = 1/2 - 1/log(1/x) on
% (0, e^{-4}) and 0 beyond is admissible with c = 1/2, its exponent is
% F(omega) = (1/2)log omega - log log omega + O(1), so e^{-F} asymp omega^{-1/2}log omega; a
% density bounded by C|x|^{-1/2} near the origin would give (phi_1 * g_eps)(0) = O(eps^{-1/2}),
% while the positive Fourier integral is at least a multiple of eps^{-1/2}log(1/eps). So the upper
% half of the comparison does not hold, and neither did SpatialLine.sato_origin_singular, admitted
% on 2026-09-15 at the ledger's paraphrase of (53.28). The librarian re-read Thm. 53.8 from the
% page image on 2026-09-19: (53.28) is
% f(x) ~ [kappa sin(k(0+)pi)/(Gamma(c)sin(c pi))] x^{c-1} K(x), with K of (53.25) and, in the
% proof of Thm. 53.6 (p. 410), the one-sided K of (53.18) said to be slowly varying at 0; the
% hypothesis (53.24) is c < infinity WITH k(0+) > 0 and k(0-) > 0. What changed here: K is defined
% once, before the three regimes, for the singular clause and the threshold clause both; the
% singular comparison function is |x|^{c-1}K(|x|); K >= 1 is stated, that being what keeps
% "unbounded" true and what cor:origin-boundedness reads; the annotation says that K is slowly
% varying and need not be bounded, with the referee's example. The Lean axiom is RESTATED AT THE
% SOURCE'S LETTER (a delegated safe-direction change, the author's decision of 2026-09-11): it
% carries K, and both it and the threshold axiom carry (53.24)'s two positivity hypotheses, which
% a symmetric law supplies with k(0+) = k(0-) = c/2. No conclusion downstream weakens: the only
% property of K that origin_boundedness spends is K >= 1 (SpatialLine.one_le_satoK, Lean core),
% which follows from k <= c on (0,infty). The last paragraph of the proof below is corrected in
% the same pass: it asserted K == 1 and L(x) = log(1/|x|) on a Cin ray, which holds only for
% tau >= 1.
% CHANGED 2026-09-19 (R181, second external review round, referee A findings 6 and 7): TWO
% REPAIRS. (a) THE POSITIVITY OF c IS A HYPOTHESIS, not an assertion. "Put c := k(0+) in (0,infty]"
% asserted c > 0, and the zero exponent is admissible with Gaussian coefficient 0 and k = 0, so
% c = 0 occurs -- and then the singular regime's range c < 1 would cover a law (delta_0) with no
% density and the clause would be false. The three Lean regimes each carry it already
% (cin_origin_singularity_unbounded has hc0 : 0 < c; the smooth one has N < c with N >= 1; the
% threshold has c = 1), as does cor:origin-boundedness, which says "finite and positive". Safe
% direction: a hypothesis added, matching the declarations. (b) THE ANNOTATION records that the
% passage from the source's asymptotic to the two-sided comparison needs the source's constant to
% be nonzero, with the computation that makes it so for symmetric data -- sin(c pi/2)/sin(c pi) =
% 1/(2 cos(c pi/2)) > 0 below the threshold, cos(c' pi/2) = 1 at it. The statement already claimed
% only the order; what was missing was the reason the order follows at all.
% CHANGED (2026-09-19, article mirror of module B, R174): the statement says, where c is
% introduced, that the three regimes concern finite c and that nothing is claimed at c = infinity
% (the register pass's item 10: the display admits c = infinity and no clause reaches it). The
% range was already in the annotation; it is now where a reader meets the hypothesis. Nothing
% else in the statement, the tag or the status changes, and no declaration is affected: the three
% Lean regimes each carry their own finiteness hypothesis already.
\begin{proposition}[the behaviour of an admissible kernel at the origin]
\label{prop:cin-origin-singularity}
\uses{thm:main-characterization,lem:selfdecomposable-exponents,prop:kernel-regularity,lem:folding-translation,lem:cin-rays}
\lean{SpatialLine.cin_origin_singularity_unbounded, SpatialLine.cin_origin_singularity_smooth,
SpatialLine.cin_origin_singularity_threshold, SpatialLine.one_le_satoK,
SpatialLine.cin_origin_singularity_cin_ray}\leanok
Let $F$ be admissible with Gaussian coefficient $0$ and folded profile $k$, let $\phi_1$ be its
kernel at canonical scale $1$, and suppose the right limit $c := k(0{+})$ is positive, so that
$c \in (0,\infty]$. The three regimes below
are stated for finite $c$, and nothing is claimed when $c = \infty$. For $0 < y < 1$ set
\[
  K(y) := \exp\Bigl[\int_y^1 \bigl(c - k(u)\bigr)\,\frac{du}{u}\Bigr],
  \qquad
  L(y) := \int_y^1 K(z)\,\frac{dz}{z},
\]
both integrals being convergent; $K \ge 1$, since $k$ is nonincreasing with right limit $c$. Then,
near the origin: if $c < 1$, $\phi_1$ is unbounded with $\phi_1(x) \asymp |x|^{c-1}K(|x|)$; if
$N < c \le N+1$ for an integer $N \ge 1$, $\phi_1$ extends continuously and has $N-1$ derivatives;
and if $c = 1$, $\phi_1$ is unbounded with $\phi_1(x) \asymp L(|x|)$. Every $\Cin$ ray has $c = 1$
with $K$ constant near the origin, so its kernel is logarithmically singular there.
\statusA\quad\ledger{A10}\quad\emph{(\textbf{Assignment.} \ledger{A10} carries the three
regimes with both comparison functions, for self-decomposable laws on $\RR$ stated there in terms
of the two-sided profile at the origin. What it does not carry is the translation into this
article's folding convention --- the source's constant is $k(0{+}) + k(0{-})$ on the two-sided
profile, which is exactly the folded $k(0{+})$ used here, and the symmetric case makes the
source's antisymmetric constant vanish. The node deliberately claims the ORDER of the
singularity and not the constant in the source's asymptotic formula; the formula itself was read
from the page image on 2026-09-19, and its constant is
$\kappa\sin(k_2(0{+})\pi)/(\Gamma(c)\sin(c\pi))$ below the threshold and $(\kappa/\pi)$ times
$\cos(c'\pi/2)$ at it, with $\kappa$ the exponential-integral expression of the source's (53.27).
\textbf{That constant has to be nonzero for the comparison to follow from the asymptotic}, since
$f \sim A g$ gives $f \asymp g$ only for $A \ne 0$, and it is nonzero here for the reason the
folding supplies (2026-09-19, R180, the referee's finding): $\kappa > 0$, and with
$k_2(0{+}) = k_2(0{-}) = c/2$ the first factor is
$\sin(c\pi/2)/\sin(c\pi) = 1/\bigl(2\cos(c\pi/2)\bigr) > 0$ for $0 < c < 1$, while at the
threshold $c' = 0$ gives $\cos(c'\pi/2) = 1$. Both are consequences of the symmetry, which is
therefore load-bearing twice: once for the antisymmetric constant and once for the nonvanishing.
The Lean specifications state the conclusion as the two-sided comparison with existential
constants $0 < c_1 \le c_2$ on a punctured neighbourhood, so they assert exactly this
consequence and not the source's asymptotic.
Range: the three regimes are stated for finite $c$; nothing is
claimed when $k(0{+}) = \infty$, which is the case for the stable profiles and for every
completely monotone profile with an unbounded Thorin measure.
\textbf{The factor $K$.} It is common to the two singular clauses --- the source's (53.25), which
enters the threshold's comparison function through (53.26) --- and it was missing from the singular
clause, and from the admitted axiom, until 2026-09-19 (the external review of module B; the marker
above). $K$ is slowly varying at the origin, and it is \emph{not} bounded there in general: the
admissible folded profile $k(x) = \tfrac12 - 1/\log(1/x)$ on $(0,e^{-4})$, $0$ beyond, has
$c = \tfrac12$ and $K(y) \asymp \log(1/y)$, and its kernel is of order $|x|^{-1/2}\log(1/|x|)$, so
no multiple of $|x|^{-1/2}$ dominates it. What the statement asserts of $K$ is $K \ge 1$, which is
elementary and machine-checked (\texttt{one\_le\_satoK}): $k$ is nonincreasing on $(0,\infty)$ with
right limit $c$, so $k \le c$ there and the integrand of $K$ is nonnegative. That bound is the only
property of $K$ the corollary below reads, and the reason "unbounded" survives the correction. The
convergence of the two integrals is elementary as well: the integrand of $K$ is nonnegative and
bounded on $[y,1]$, the profile being monotone there, and $K$ is nonincreasing, so
$K(z)/z \le K(y)/y$ on $[y,1]$; both steps
are machine-checked (\texttt{integrableOn\_satoIntegrand} for every $c$ on the two-sided profile,
with \texttt{integrableOn\_thresholdIntegrand} and \texttt{integrableOn\_thresholdComparison} at
the threshold). The last sentence, that every
$\Cin$ ray sits at the threshold with $K$ constant near the origin, is $[T]$ and is
machine-checked (2026-09-09); it rests on Lean core and on no clause of \ledger{A10}. The three
regimes rest on the entry and on nothing else, and \textbf{the entry is admitted} (2026-09-15,
the author's decision): its three regimes are the Lean axioms \texttt{sato\_origin\_singular},
\texttt{sato\_origin\_smooth} and \texttt{sato\_origin\_threshold}, one per regime, and the three
statements above are derived from them on Lean core. The admission waited through six waves of
the proving campaign for two reasons, and both were removed the day before it was taken: nothing
spent the regimes --- \ref{cor:origin-boundedness} now spends all three --- and admitting them in
this article's \emph{folded} form would have put the folding translation, which \ledger{A10} says
it does not carry, behind an admitted name. So the three names are stated \textbf{at the source's
own letter}, over a two-sided profile and its two one-sided limits, and the translation is
\ref{lem:folding-translation} beside them. What they carry beyond their pages is four things,
the same four for each, and the canonical list is at the head of
\texttt{blueprint/trust-boundary.txt}: the hypothesis is this article's cosine-transform
specification of a symmetric law rather than the source's self-decomposable law, which is a
passage through \ledger{A3}; the conclusion is the two-sided comparison above, with its constants
and its neighbourhood existential, which is \emph{less} than the source's asymptotic with its
constant and cosine factor but needs that constant to be nonzero, as the paragraph above records;
the datum is a monotone two-sided profile rather than the source's one-sided-continuous version of
it, a modification on a countable set that changes nothing the theorems speak of
(\ref{lem:folding-translation}'s annotation); and the folding convention, which is no longer among
them. So this
node is \leanok{} in the sense \ref{prop:kernel-regularity} is: what is machine-checked is the
passage from the source's sentences to these, and the analysis itself remains cited.
\textbf{The singular name was restated on 2026-09-19} (R171), when the referee's example showed
that what it asserted was false: it carries the factor $K$ now, and it and the threshold name carry
the source's own hypothesis that both one-sided limits are positive, which a symmetric law supplies
with $k_2(0{+}) = k_2(0{-}) = c/2$. No conclusion downstream became weaker, the corollary reading
$K \ge 1$ alone. Their
hypotheses were checked at the edges in wave 6 (2026-09-10) and again before the admission:
$c = 0$ is excluded from all three, and must be, since with $a = 0$ it forces $k \equiv 0$ and a
law without a density; the Gaussian is excluded by $a = 0$; $c = \infty$ is excluded by the shape
of the hypothesis; and none of the three is vacuous, the threshold regime being met by the $\Cin$
rays, the singular one by $c\,\mathbf 1_{(0,\tau)}$ with $c < 1$ and the smooth one by the
\Matern{} profile at $\gamma = 1$, where $c = 2$. The node never depended on
\ref{lem:cin-delay-equation}: the three regimes quantify over an arbitrary admissible $F$ with
Gaussian coefficient $0$ and speak only of its kernel near the origin --- what did depend on the
delay equation was the propagation, which is \ref{rem:cin-propagation}.)}
\end{proposition}

% ADDED (2026-09-15, module B step 2, R162): a proof environment, for the reason R157 added one
% at prop:kernel-regularity. An [A] node that prints nothing between its cited facts and its
% conclusion leaves this article's own step in the annotation alone, and here that step is the
% passage between the two conventions -- which is exactly what the admission was made to avoid
% taking on trust. The proof is the one the Lean takes: lem:folding-translation, then the source's
% three sentences read through it, then the Cin-ray computation.
\begin{proof}
\leanok
The three regimes are the cited interface, and what has to be said here is how the source's
statement becomes this one. The sources state the behaviour at the origin of a self-decomposable
law with no Gaussian part in terms of its \emph{two-sided} profile $k_2$ on $\RR \setminus \{0\}$
and the two constants $k_2(0{+}) + k_2(0{-})$ and $k_2(0{+}) - k_2(0{-})$: the first is their
$c$, the second their antisymmetric constant. By \ref{lem:folding-translation} the law in
question is theirs. Its two-sided datum is $k_2(x) = \tfrac12 k(|x|)$, which is a profile of the
kind the sources require; its exponent is the folded exponent, so the two descriptions specify
the same law; the first constant is $k(0{+}) = c$; and the second vanishes, the law being
symmetric. Three hypotheses of the sources are met and not visible in that sentence, and
\ref{lem:folding-translation}'s annotation says why: with $c$ finite the two-sided profile has
finite variation near the origin, $\int_{|x|<1}k_2 \le c$; the law has no drift, being specified
through its cosine transform, and no Gaussian part, since $a = 0$; and the one-sided continuity by
which the sources normalise a $k$-function is a modification on a countable set, under which the
\Levy{} measure, the exponent, the two limits at the origin and both comparison functions are
unchanged. Reading the three regimes through that dictionary gives the three statements
displayed above, with $k_2(u) + k_2(-u)$ becoming $k(u)$ in the integrand of $K$, which is the
factor both singular clauses carry. The passage from the source's asymptotic to the two-sided
comparison uses that the constant in the source's formula does not vanish, which the symmetry
supplies; the annotation computes it. That passage is machine-checked; the regimes themselves are
cited, and the order of the singularity is all that is claimed of them.

The last sentence is this article's own. The $\Cin$ ray of step $\tau$ has profile
$\mathbf 1_{(0,\tau)}$ (\ref{lem:cin-rays}(1)), whose right limit at the origin is $1$, so
$c = 1$ and the threshold regime applies; and on $(0,\tau)$ the integrand $(1 - k(u))u^{-1}$ of
$K$ vanishes, which is the constancy of an indicator on its own interval and is what is
machine-checked, on Lean core. So $K$ is constant on $(0,\min(\tau,1))$, with the value
$\max(1,1/\tau)$: for $\tau \ge 1$ nothing is integrated, and for $\tau < 1$ the window $(\tau,1)$
contributes $\int_\tau^1 u^{-1}du = \log(1/\tau)$. Hence $L$ is a constant multiple of
$\log(1/|x|)$ plus a constant near the origin --- $L(x) = \log(1/|x|)$ when $\tau \ge 1$ and
$L(x) = \tau^{-1}\log(\tau/|x|) + \tau^{-1} - 1$ when $\tau < 1$ --- and in both cases the
logarithmic order is all the statement claims.
\end{proof}

% CHANGED (review 2026-09-09, Q8): the propagation of the singularity to the points +-n tau was
% MOVED OUT of the proposition's statement into this remark. It is [T] content, its induction is
% written down nowhere, and it sat inside an [A] statement -- the worst place for an unproved
% claim, since a reader takes the ledger to cover the whole node. Nothing scheduled consumes it
% (its only consumer is prop:jet-regularity(4), Chapter 15, unscheduled), so it is described
% rather than asserted; should it be needed, the induction has to be written and typed first.
\begin{remark}[the singularity propagates outward]\label{rem:cin-propagation}
For a $\Cin$ ray of step $\tau$ the delay equation of \ref{lem:cin-delay-equation} carries the
singularity outward along the lattice $\tau\ZZ$, one order of smoothness gained at each step:
the derivative of the kernel inherits a logarithmic singularity at $\pm\tau$, so the kernel
itself is continuous there with a kink of order $|x \mp \tau|\log|x \mp \tau|$, and so on to
$\pm2\tau$ and beyond. This is a description of the shape of the kernel and not a claim of the
proposition: the induction it summarises is not carried out here, and no node depends on it.
\end{remark}

The $\Cin$ kernels are therefore logarithmically singular at the origin, like the
variance-gamma member of smoothness index $0$, whose density is a Bessel $K_0$ and which sits
at the same threshold. That coincidence is not one: both are instances of a single boundary,
and where it falls is decided by the profile alone.

% NEW NODE (2026-09-15, module B step 2, R159; the postdoc's P-0008). Three things in the
% article said this and were never put side by side: the external review's finding C, that the
% closed form of prop:matern-density is a \Matern{} covariance only for gamma > 1/2 and is K_0 at
% gamma = 1/2; ledger A10's three regimes, whose boundary is c = 1; and the \Matern{} profile's
% own value 2 gamma at the origin. The corollary is what they say together, and it claims
% strictly less than prop:cin-origin-singularity, being a two-line deduction from the three
% regimes. It gives A10 what wave 6 found it lacked, a consumer reading all three regimes, and
% it turns a naming caveat about one family into a statement about every admissible kernel.
% The \uses edge to prop:matern-exponent is FORWARD, into Chapter 10; it is deliberate, and it is
% the same deliberate forward reference prop:two-members makes to the same node for the same
% reason -- the corner computation belongs to Chapter 10 and the statement that needs it does
% not. There is no cycle: prop:matern-exponent reaches Chapters 2 and 7 only.
% CHANGED (2026-09-15, module B step 2, R162): the node is \leanok and its statement is NARROWED
% by one half-sentence. THE STATUS: ledger A10 is admitted (the author's decision), so clause (1)
% is SpatialLine.origin_boundedness -- proved, spending all three names -- where it was
% Skeleton.origin_boundedness carrying sorry, which is deleted. THE NARROWING: clause (2)'s
% closing "where prop:matern-density gives it as the Bessel K_0" is dropped from the STATEMENT
% and moved into the annotation. It is a pointer to a node that is itself notready, and a \leanok
% statement should not assert, even by attribution, what no declaration of its own tag proves;
% the sentence survives as a cross-reference where it belongs. THE TAG gains
% SpatialLine.matern_origin_boundedness, the bundle that reads clause (1) at the \Matern{} member
% and concludes what clause (2) claims of it -- boundedness exactly for gamma > 1/2 -- so that
% each clause has a declaration proving it rather than leaving the modus ponens to the reader
% (the fidelity-review rule on \lean tags). Statement otherwise unchanged.
\begin{corollary}[the boundedness threshold]\label{cor:origin-boundedness}
\uses{prop:cin-origin-singularity,lem:folding-translation,lem:cin-rays,prop:matern-exponent}
\lean{SpatialLine.origin_boundedness, SpatialLine.matern_origin_boundedness,
SpatialLine.matern_origin_constant,
SpatialLine.cin_origin_singularity_cin_ray}\leanok
Let $F$ be admissible with Gaussian coefficient $0$ and folded profile $k$, let $\phi_1$ be its
kernel at canonical scale $1$, and suppose $c := k(0{+})$ is finite and positive.
\begin{enumerate}
\item \emph{(The threshold.)} If $c > 1$, then $\phi_1$ has a representative continuous on the
whole line. If $c \le 1$, then for every representative $p$ of $\phi_1$, every $\delta > 0$ and
every $M$ there is an $x$ with $0 < |x| < \delta$ and $p(x) > M$; no representative of $\phi_1$
is bounded near the origin. So $\phi_1$ is bounded at the origin exactly when $c > 1$.
\item \emph{(The named members.)} Every $\Cin$ ray has $c = 1$ and sits at the threshold. The
\Matern{} member of index $\gamma$ and range $\theta$ has $c = 2\gamma$, so its kernel is
bounded at the origin exactly for $\gamma > \tfrac12$, and at $\gamma = \tfrac12$ it sits at the
threshold.
\end{enumerate}
\statusA\quad\ledger{A10}\quad\emph{(\textbf{Assignment.} Clause (1) is a deduction from the
three regimes of \ref{prop:cin-origin-singularity} and carries no content beyond them, so it
rests on \ledger{A10} exactly as that node does and on the same finite range: nothing is claimed
when $k(0{+}) = \infty$. The deduction reads all three regimes, which is what makes this the
consumer the entry lacked, and what the entry's admission (2026-09-15) bought. Two steps of it
are this article's own and are machine-checked: that both singular comparison functions diverge,
which is one bound on the slowly varying factor, $K \ge 1$, read off the profile's own
monotonicity and not off \ledger{A16}'s asymptotic (below the threshold it leaves
$|x|^{c-1}K(|x|) \ge |x|^{c-1}$, and at the threshold $L(x) \ge \log(1/|x|)$); and the passage
from \emph{one} unbounded representative to
\emph{every} one, which compares measures rather than functions. Clause (2) is $[T]$ in both
halves and machine-checked: that every $\Cin$ ray has $c = 1$ is
\texttt{cin\_origin\_singularity\_cin\_ray}, proved in 2026-09-09, and that the \Matern{} profile
has $k(0{+}) = 2\gamma$, with the three regimes reading $\gamma < \tfrac12$,
$\gamma = \tfrac12$ and $\gamma > \tfrac12$, is \texttt{matern\_origin\_constant}; both rest on
Lean core and on no clause of \ledger{A10}, and the member's own boundedness statement is clause
(1) read at $c = 2\gamma$ (\texttt{matern\_origin\_boundedness}). What clause (2) adds to clause
(1) is therefore only arithmetic, and the boundary it names is the one the external review found
from the other side, at the closed form of \ref{prop:matern-density}: that closed form is a
\Matern{} covariance only for $\gamma > \tfrac12$ and is the Bessel $K_0$ at $\gamma = \tfrac12$.
That sentence stood in clause (2) until the admission and is here instead: it is a pointer to
\ref{prop:matern-density}, whose closed form is not machine-checked, and a statement of this node
should assert only what this node's own declarations prove.
\textbf{Not claimed}: the constant in the comparison, as at \ref{prop:cin-origin-singularity};
the rate of the divergence, at the threshold or below it. That rate is not logarithmic in general:
$L$ is of logarithmic order exactly where $K$ is constant near the origin, which is the case for
the $\Cin$ rays and for the \Matern{} member at $\gamma = \tfrac12$ but not for every profile at
the threshold --- with $K \asymp \log(1/x)$ one gets $L \asymp \log^2(1/x)$ --- and the rate for
the named members is \ledger{A16}'s asymptotic rather than anything proved here. And any statement
in dimension greater than one.)}
\end{corollary}

% CHANGED (2026-09-15, module B step 2, R162): the proof carries \leanok, and its measure
% comparison is rewritten to the CHECKED ROUTE (the author's decision of 2026-09-09). The printed
% argument integrated the regime's lower bound over (0,epsilon) to get c_1 epsilon^c / c; the Lean
% proof does not integrate it at all, but replaces it on that interval by its smallest value,
% which is the same comparison with one step fewer -- and the same at the threshold with
% log(1/epsilon). Nothing else of the proof changed.
\begin{proof}
\leanok
Both halves of clause (1) come from \ref{prop:cin-origin-singularity}. If $c > 1$, choose the
integer $N \ge 1$ with $N < c \le N+1$ --- which exists for every finite $c > 1$, namely
$N = \lceil c\rceil - 1$ --- and the smooth regime gives a representative of $\phi_1$ that is
$C^{N-1}$ on the whole line, in particular continuous. Below the threshold and at it, both
comparison functions are handled by one bound on the slowly varying factor: $k$ is nonincreasing
on $(0,\infty)$ with right limit $c$, so $k \le c$ there, the integrand $(c-k(u))u^{-1}$ of $K$ is
nonnegative, and $K \ge 1$ on $(0,1)$. If $c < 1$ the singular regime gives a
representative $p_0$ and constants $c_1 > 0$, $\delta_0 > 0$ with
$p_0(x) \ge c_1|x|^{c-1}K(|x|) \ge c_1|x|^{c-1}$ for $0 < |x| < \delta_0$; if $c = 1$ the
threshold regime gives the same with $c_1 L(|x|)$ in place of $c_1|x|^{c-1}K(|x|)$, and
$L(x) \ge \int_{|x|}^1 z^{-1}dz = \log(1/|x|)$ by the same bound. In both cases
the lower bound tends to $+\infty$ as $x \to 0$.

That carries to every representative, and not only to $p_0$, by comparing measures rather than
functions. Suppose some representative $p$ satisfied $p \le M$ on the punctured
$\delta$-neighbourhood of the origin, and let $0 < \varepsilon \le \min(\delta,\delta_0)$. On
$(0,\varepsilon)$ the regime's lower bound is at least its value at $\varepsilon$, both
$|x|^{c-1}$ and $L$ being nonincreasing in $|x|$: below the threshold
$p_0 \ge c_1\varepsilon^{c-1}$ there, and at it $p_0 \ge c_1\log(1/\varepsilon)$. So
$\phi_1\bigl((0,\varepsilon)\bigr) \ge c_1\varepsilon^{c-1}\cdot\varepsilon$ in the first case
and $\ge c_1\log(1/\varepsilon)\cdot\varepsilon$ in the second, while the assumed bound on $p$
gives $\phi_1\bigl((0,\varepsilon)\bigr) \le M\varepsilon$. Dividing by $\varepsilon$ leaves
$c_1\varepsilon^{c-1} \le M$, respectively $c_1\log(1/\varepsilon) \le M$, and both left-hand
sides tend to $+\infty$ as $\varepsilon$ decreases. Choosing $\varepsilon$ small enough
contradicts the assumption. Neither step needs $p$ to be measurable: a density integrates over a
Borel set whatever representative names it, and the two bounds are bounds on that integral.

Clause (2) is the profile computation in each case. The $\Cin$ ray of step $\tau$ has profile
$\mathbf 1_{(0,\tau)}$, whose right limit at the origin is $1$ (\ref{lem:cin-rays}(1)); the
\Matern{} profile at range $\theta$ is $2\gamma e^{-x/\theta}$ (\ref{prop:matern-exponent}(1)),
whose right limit is $2\gamma$. Reading $c = 2\gamma$ through clause (1) puts the threshold at
$\gamma = \tfrac12$.
\end{proof}

% NEW NODE 2026-09-19 (R189; the author's decision D1, a widening authorized at the second
% external review round of module B). The corollary above assumes c finite, so at c = infinity
% the article claimed nothing -- and that is the case for the stable members, whose row in the
% article's comparison table says "bounded at 0: yes". Ledger A10's smooth regime is Sato
% Thm. 28.4, and its clause (ii), read from the page image by the librarian on 2026-09-19 and
% transcribed into blueprint/AXIOMS-verbatim.md the same day, covers c = infinity in one
% sentence: "If c = infinity, then mu has a C^infinity density on R". This node is that clause
% read through the folding translation. It is PROSE-ONLY: no Lean is scheduled, no name is
% admitted, blueprint/trust-boundary.txt is unchanged, and the three admitted A10 names cover
% the finite regimes alone. The node is \notready in this article's sense -- stated, cited, not
% machine-checked -- as lem:thorin-ggc is.
% Note on the \uses edge: the node does NOT use lem:folding-translation, whose hypothesis is
% c < infinity; the two lines of translation it needs are written out in the proof, and the
% lemma's own annotation records that its finiteness hypothesis is used for the two limits alone.
\begin{corollary}[an unbounded profile makes the kernel smooth]\label{cor:origin-smooth}
\uses{thm:main-characterization,lem:profile-integrability,prop:stable-family}
\notready
Let $F$ be admissible with Gaussian coefficient $0$ and folded profile $k$, and let $\phi_1$ be
its kernel at canonical scale $1$. If $k(0{+}) = \infty$, then $\phi_1$ has a representative of
class $C^\infty$ on $\RR$; in particular that representative is bounded on every bounded set, so
$\phi_1$ is bounded at the origin. The symmetric stable members of index $0 < \alpha < 2$ are of
this kind: their folded profile is a positive multiple of $x^{-\alpha}$
(\ref{prop:stable-family}), whose right limit at the origin is $+\infty$.
\statusA\quad\ledger{A10}\quad\emph{(\textbf{Assignment.} \ledger{A10} carries the conclusion:
for a self-decomposable law on $\RR$ presented through a two-sided profile with
$k_2(0{+}) + k_2(0{-}) = \infty$, the law has a $C^\infty$ density. What it does not carry is
the passage into this article's folded convention, which is two lines of the proof below and is
$[T]$; the normalisation of the source's $k$-function to its one-sided-continuous version, which
is the same modification on a countable set that \ref{lem:folding-translation} discusses and that
the three admitted names of \ledger{A10} carry beyond their pages
(\texttt{blueprint/trust-boundary.txt}); and the two consequences, that a $C^\infty$ function is
bounded on every compact set and that the stable profile is unbounded at the origin, both $[T]$.
\textbf{Not machine-checked, and no name is admitted.} The three Lean axioms behind
\ref{prop:cin-origin-singularity} are the regimes at \emph{finite} $c$; the clause at
$c = \infty$ is cited here and nowhere else in the library, so the trust boundary is unchanged
and the node is \texttt{notready}. \textbf{Not claimed}: any rate, any bound on the density or
its derivatives, and anything about the Student-t members. Their profile is unbounded at the
origin too, but this article does not prove it --- what it has of that family's profile is the
mixture \ref{eq:bridge-profile} of a causal catalogue it does not compute --- and nothing is
lost, the Student-t kernel being bounded at the origin by its own explicit density,
\ref{prop:student-t}(1).
\textbf{The dichotomy.} With the Gaussian coefficient $0$ and $F \not\equiv 0$ the right limit
$c = k(0{+})$ lies in $(0,\infty]$, since a nonincreasing $k$ with $k(0{+}) = 0$ vanishes
identically. So this node and \ref{cor:origin-boundedness} together cover the whole range and say
one thing: \emph{an admissible kernel with no Gaussian part is bounded at the origin exactly when
its profile starts above $1$.})}
\end{corollary}

\begin{proof}
The cited theorem is stated for a self-decomposable law on $\RR$ with no Gaussian part, presented
through a two-sided profile $k_2$ on $\RR \setminus \{0\}$ --- nonnegative, nonincreasing on
$(0,\infty)$ and nondecreasing on $(-\infty,0)$ --- and reads its conclusion off the constant
$c = k_2(0{+}) + k_2(0{-})$; at $c = \infty$ that conclusion is a $C^\infty$ density on the line.
That presentation is itself the hypothesis: the source says of it that it is the general form of a
purely non-Gaussian self-decomposable law on $\RR$, so exhibiting a law in it is exhibiting it as
one.
Put $k_2(x) := \tfrac12 k(|x|)$. The structural conditions and the exponent identity are the first
two paragraphs of the proof of \ref{lem:folding-translation}, which read only the monotonicity of
$k$ and the two integrability conditions of \ref{lem:profile-integrability}: that lemma's
hypothesis that $k(0{+})$ is finite is used for its two one-sided limits alone, and is not
available here. So $k_2$ is a two-sided profile of the required kind and its \Levy{} measure
$k_2(x)|x|^{-1}dx$ specifies the same symmetric law as the pair $(0,k)$, whose kernel at canonical
scale $1$ is $\phi_1$. Both one-sided limits of $k_2$ at the origin are $+\infty$: on $(0,\infty)$
it is $\tfrac12 k$, whose right limit is $+\infty$ by hypothesis, and on $(-\infty,0)$ it is
$x \mapsto \tfrac12 k(-x)$, where $x \mapsto -x$ carries the left neighbourhood filter at the
origin to the right one. Hence the source's constant is $c = \infty$. One hypothesis of the source
is not displayed by that dictionary, and it is the one \ref{lem:folding-translation}'s annotation
records: the source normalises a $k$-function to the version left-continuous on $(0,\infty)$ and
right-continuous on $(-\infty,0)$, and $k_2$ need not be either. Passing to that modification
changes $k_2$ on a countable set, hence changes neither the \Levy{} measure nor the exponent nor
the law, and leaves both one-sided limits infinite; so the theorem applies to the modification and
concludes about the same law.

The two consequences are elementary. A function of class $C^\infty$ on $\RR$ is continuous, hence
bounded on every compact set, so the representative the theorem produces is bounded on a
neighbourhood of the origin. And the symmetric stable member of index $\alpha$ has folded profile
$C_\alpha^{-1}x^{-\alpha}$ with $C_\alpha \in (0,\infty)$ (\ref{prop:stable-family}), which tends
to $+\infty$ as $x$ decreases to $0$; so the hypothesis holds for every member of that family.
\end{proof}
